% =====================================================================
% Observation Space Domain Randomization using Colored Noise (Overleaf-compatible)
%
% Required Packages (Please add to your main document preamble):
%   \usepackage{amsmath}
%   \usepackage{amssymb}
%   \usepackage{bm}
%   \usepackage{booktabs}
% =====================================================================

\subsection{Domain Randomization in Observation Space via Colored Noise}
\label{sec:obs_pink_noise}

\subsubsection{Motivation}

One of the primary challenges in Sim-to-Real transfer is the discrepancy in sensor characteristics between the simulator and the real robot (commonly referred to as the Reality Gap).
In practice, measurement errors from encoders and sensors on real physical systems are rarely independent and identically distributed (i.i.d.); rather, they exhibit strong temporal correlations between consecutive time steps~\cite{eberhard2023pink}.
For example, joint encoder errors on servo motors tend to depend on the motor's current rotational direction or velocity, leading to persistent, slowly-varying deviations over time.

Conventional domain randomization methods typically corrupt observations by adding independent white noise (e.g., uniform or Gaussian noise) at each control step.
However, because white noise possesses a flat power spectrum containing all frequency components equally, it fails to capture the low-frequency-dominated error structures characteristic of physical sensors.
As a consequence, policies trained with white noise often lack robustness against persistent sensor biases when deployed on real systems, resulting in degraded Sim-to-Real performance.

In this work, we extend the concept of colored noise introduced by Eberhard~et~al.~\cite{eberhard2023pink}---which was originally proposed as an action noise exploration strategy---to the observation space as a domain randomization technique.
Specifically, we inject temporally correlated noise whose power spectral density follows a power-law relationship $S(f) \propto 1/f^{\beta}$ (with $\beta = 1$, representing pink noise) into the robot's joint positions and velocities.

\subsubsection{Observation Space Definition}

The structure of the observation space in this task is detailed in Table~\ref{tab:obs}.
At each step, the policy receives a concatenated observation vector $\boldsymbol{o}_t \in \mathbb{R}^{n_\mathrm{obs}}$ representing the robot's state and task commands, and outputs joint position targets $\boldsymbol{a}_t \in \mathbb{R}^{n_\mathrm{act}}$.

\begin{table}[htbp]
  \centering
  \caption{Observation space configuration}
  \label{tab:obs}
  \begin{tabular}{llcc}
    \toprule
    Observation Item & Description & Dimension & Noise Injection \\
    \midrule
    Joint position $\boldsymbol{q}_t$      & Relative angles of each joint & 6  & \checkmark \\
    Joint velocity $\dot{\boldsymbol{q}}_t$ & Relative angular velocities of each joint & 6  & \checkmark \\
    Target end-effector pose $\boldsymbol{p}_t^{\mathrm{cmd}}$ & Position $(x,y,z)$ + Quaternion $(q_w,q_x,q_y,q_z)$ & 7  & -- \\
    Previous action $\boldsymbol{a}_{t-1}$ & Policy output from the previous step & 6 & -- \\
    \midrule
    Total & & 25 & \\
    \bottomrule
  \end{tabular}
\end{table}

We apply pink noise only to the joint positions (6D) and joint velocities (6D), totaling 12 dimensions representing the physical state of the robot.
We do not inject noise into the target command or the previous action, as these signals are generated cleanly within the control loop or the policy network.

\subsubsection{Mathematical Formulation of Colored Noise}

Colored noise refers to a stochastic process whose power spectral density (PSD), $S(f)$, scales with frequency $f$ according to:

\begin{equation}
  S(f) \propto \frac{1}{f^{\beta}}
  \label{eq:psd}
\end{equation}

where $\beta$ denotes the spectral exponent.
Special cases include white noise when $\beta = 0$ (equal energy across all frequencies), pink noise when $\beta = 1$ ($1/f$ noise), and brown/red noise when $\beta = 2$.

Compared to white noise, pink noise ($\beta = 1$) highlights lower frequency components while dampening high-frequency oscillations.
This spectral shape naturally models the slowly-varying bias drift and systematic measurement errors observed in real-world physical sensors.

To generate the pink noise sequences, we employ the Fast Fourier Transform (FFT) based algorithm proposed by Timmer \& Koenig~\cite{timmer1995}.
A $d$-dimensional pink noise sequence of length $T$, denoted as $\boldsymbol{\eta} \in \mathbb{R}^{d \times T}$, is generated through the following steps:

\begin{enumerate}
  \item Compute the discrete frequency bins $f_k = k / T$ for $k = 0, 1, \ldots, \lfloor T/2 \rfloor$.
  \item Construct amplitude scaling factors $s_k = f_k^{-\beta/2}$.
        To prevent low-frequency divergence in finite-length sequences, a low-frequency cutoff is defined at $f_{\min} = 1/T$, and the value $s_{k_{\min}}$ is assigned to all frequency bins below this threshold.
  \item Generate the real part $\mathrm{Re}(X_k)$ and imaginary part $\mathrm{Im}(X_k)$ of the Fourier coefficients independently from $\mathcal{N}(0, s_k^2)$.
        The imaginary component of the DC term ($k=0$) is set to zero. If the sequence length $T$ is even, the Nyquist frequency component is also constrained to be real.
        These components are multiplied by $\sqrt{2}$ to correct their amplitudes.
  \item Perform an Inverse Discrete Fourier Transform (IDFT) to yield the time-domain series, and normalize the output to unit variance by dividing by the theoretical standard deviation $\sigma_{\mathrm{th}}$:
        \begin{equation}
          \sigma_{\mathrm{th}} = \frac{2}{T} \sqrt{\sum_{k=1}^{\lfloor T/2 \rfloor} w_k^2}
          \label{eq:sigma_th}
        \end{equation}
        where $w_k = s_k$ for $k < \lfloor T/2 \rfloor$, and the final bin is adjusted as $w_{\lfloor T/2 \rfloor} = s_{\lfloor T/2 \rfloor} \cdot (1 + (T \bmod 2)) / 2$.
\end{enumerate}

\subsubsection{Observation Noise Injection Model}

Using the generated unit-variance pink noise process, we corrupt the joint observations at step $t$ as follows:

For joint positions:
\begin{equation{eq:obs_noise_pos}}
  \tilde{\boldsymbol{q}}_t = \boldsymbol{q}_t + \sigma_{\mathrm{obs}} \, \boldsymbol{\eta}_t^{(q)}
\end{equation}

For joint velocities:
\begin{equation{eq:obs_noise_vel}}
  \tilde{\dot{\boldsymbol{q}}}_t = \dot{\boldsymbol{q}}_t + \sigma_{\mathrm{obs}} \, \boldsymbol{\eta}_t^{(\dot{q})}
\end{equation}

where $\boldsymbol{\eta}_t^{(q)} \in \mathbb{R}^{6}$ and $\boldsymbol{\eta}_t^{(\dot{q})} \in \mathbb{R}^{6}$ represent independent pink noise samples for position and velocity, respectively, and $\sigma_{\mathrm{obs}}$ represents the noise scaling parameter.
Each parallel environment maintains its own independent pink noise generator, ensuring that noise sequences are uncorrelated across environments.

Consequently, the policy inputs at time step $t$ are constructed as:
\begin{equation}
  \tilde{\boldsymbol{o}}_t =
  \begin{bmatrix}
    \tilde{\boldsymbol{q}}_t \\
    \tilde{\dot{\boldsymbol{q}}}_t \\
    \boldsymbol{p}_t^{\mathrm{cmd}} \\
    \boldsymbol{a}_{t-1}
  \end{bmatrix}
  \in \mathbb{R}^{25}
  \label{eq:corrupted_obs}
\end{equation}

\subsubsection{Buffer-based Noise Process Management}

Because the FFT-based generation of pink noise is computationally expensive, generating it online at every step would severely bottleneck GPU-accelerated parallel simulation.
To circumvent this, we implement a buffer-based approach.

Prior to training, pink noise sequences spanning the maximum episode length $T_{\mathrm{ep}}$ are generated in batch and stored in memory.
During step transitions, we index sequentially through the buffer.
When a specific environment is reset or reaches the end of the buffer, only the noise buffer of that environment is regenerated with a new sequence.
This design maintains the temporal correlation structure of the noise while keeping the step-wise computation cost at $O(1)$.

\subsubsection{Evaluation Scenarios}

To isolate and evaluate the impact of colored observation noise on Sim-to-Real policy transfer, we construct four distinct training scenarios (Table~\ref{tab:noise_conditions}).

\begin{table}[htbp]
  \centering
  \caption{Comparison of noise injection conditions}
  \label{tab:noise_conditions}
  \begin{tabular}{lcc}
    \toprule
    Scenario Name & Action Noise & Observation Noise \\
    \midrule
    No-Noise (Baseline) & -- & -- \\
    Action-Noise        & Pink & -- \\
    \textbf{Obs-Noise (Ours)} & -- & \textbf{Pink} \\
    Both-Noise          & Pink & Pink \\
    \bottomrule
  \end{tabular}
\end{table}

In the Obs-Noise scenario (\texttt{Isaac-SO-ARM101-Reach-Sim2Real-Obs-Noise-v0}), action perturbations are disabled, and pink noise is injected solely into the joint positions and velocities.
This configuration allows us to isolate the effects of sensor noise simulation from command perturbations and examine its individual contribution to transfer performance.

Note that across all four conditions, robot link mass randomization ($\pm 30\%$ scaling) remains enabled as a baseline domain randomization parameter.

\subsubsection{Hyperparameter Settings}

The primary hyperparameters utilized in our implementation are summarized in Table~\ref{tab:hyperparams}.

\begin{table}[htbp]
  \centering
  \caption{Pink noise hyperparameters}
  \label{tab:hyperparams}
  \begin{tabular}{lll}
    \toprule
    Parameter & Symbol & Value \\
    \midrule
    Spectral exponent & $\beta$ & $1$ (Pink noise) \\
    Observation noise scale & $\sigma_{\mathrm{obs}}$ & $0.02$ \\
    Action noise scale (for Both-Noise) & $\sigma_{\mathrm{act}}$ & $0.03$ \\
    Buffer length & $T$ & $1000$ steps \\
    Low-frequency cutoff & $f_{\min}$ & $1/T$ \\
    \bottomrule
  \end{tabular}
\end{table}

The observation noise scale $\sigma_{\mathrm{obs}} = 0.02$ was selected based on the typical error specifications (tens of milliradians) of physical servo encoders.
This amplitude effectively models real-world sensor noise without destabilizing the policy optimization process.

% =====================================================================
% BibTeX References
% =====================================================================
% @inproceedings{eberhard2023pink,
%   title     = {Pink Noise Is All You Need: Colored Noise Exploration in Deep Reinforcement Learning},
%   author    = {Eberhard, Onno and Hollenstein, Jakob and Pinneri, Cristina and Martius, Georg},
%   booktitle = {Proceedings of the Eleventh International Conference on Learning Representations (ICLR 2023)},
%   month     = may,
%   year      = {2023},
%   url       = {https://openreview.net/forum?id=hQ9V5QN27eS},
% }
% @article{timmer1995,
%   title   = {On generating power law noise},
%   author  = {Timmer, J. and Koenig, M.},
%   journal = {Astronomy and Astrophysics},
%   volume  = {300},
%   pages   = {707--710},
%   year    = {1995},
% }
